\subsection{Работа с сетевыми устройствами}

К функционалу работы с сетевыми устройствами можно отнести следующие компоненты:

\begin{itemize}
    \item поиск сетевых устройств;
    \item получение списка сетевых шаров;
    \item поиск файлов и шаров на сетевых устройствах;
\end{itemize}

\subsubsection{Поиск сетевых устройств}

Поиск сетевых устройств начинается с инициализации множества для хранения уникальных устройств и вектора для результата.

Алгоритм использует несколько методов для максимально полного обнаружения сетевых устройств. Первый метод -- получение сетевых дисков, сопоставленных буквам (Z:, Y: и т.д.), которые добавляются в множество устройств.

\begin{lstlisting}[style=CodeListing]
std::vector<std::wstring> networkDrives = EnumerateNetworkDrives();
for (const auto& dev : networkDrives) {
    if (!dev.empty()) {
        devicesSet.insert(dev);
    }
}
\end{lstlisting}

Второй метод -- рекурсивное перечисление всех сетевых ресурсов, начиная с корня сети, что позволяет найти сетевые компьютеры и шары аналогично представлению "Сеть" в проводнике Windows.

\begin{lstlisting}[style=CodeListing]
NETRESOURCEW* pNetResource = nullptr;
EnumerateNetworkResourcesRecursive(pNetResource, devicesSet, 0);
\end{lstlisting}

Третий метод -- перечисление рабочих групп и доменов с их серверами через установку типа отображения RESOURCEDISPLAYTYPE\_DOMAIN.

\begin{lstlisting}[style=CodeListing]
NETRESOURCEW nr{};
nr.dwScope = RESOURCE_GLOBALNET;
nr.dwType = RESOURCETYPE_ANY;
nr.dwUsage = RESOURCEUSAGE_CONTAINER;
nr.dwDisplayType = RESOURCEDISPLAYTYPE_DOMAIN;
EnumerateNetworkResourcesRecursive(&nr, devicesSet, 0);
\end{lstlisting}

Четвёртый метод -- прямое перечисление серверов через установку типа отображения RESOURCEDISPLAYTYPE\_SERVER.

\begin{lstlisting}[style=CodeListing]
nr.dwScope = RESOURCE_GLOBALNET;
nr.dwType = RESOURCETYPE_ANY;
nr.dwUsage = RESOURCEUSAGE_ALL;
nr.dwDisplayType = RESOURCEDISPLAYTYPE_SERVER;
EnumerateNetworkResourcesRecursive(&nr, devicesSet, 0);
\end{lstlisting}

Пятый метод -- перечисление всех дисковых шаров через установку типа ресурса RESOURCETYPE\_DISK.

\begin{lstlisting}[style=CodeListing]
nr.dwScope = RESOURCE_GLOBALNET;
nr.dwType = RESOURCETYPE_DISK;
nr.dwUsage = RESOURCEUSAGE_ALL;
nr.dwDisplayType = RESOURCEDISPLAYTYPE_GENERIC;
EnumerateNetworkResourcesRecursive(&nr, devicesSet, 0);
\end{lstlisting}

Шестой метод -- перечисление из корня сетевого окружения с использованием флагов RESOURCEUSAGE\_CONTAINER и RESOURCEUSAGE\_CONNECTABLE. При успешном открытии перечисления выполняется перебор всех найденных ресурсов: проверяется формат имени, извлекаются имена серверов и добавляются в множество устройств, для контейнеров выполняется рекурсивное перечисление.

\begin{lstlisting}[style=CodeListing]
HANDLE hEnum = nullptr;
if (WNetOpenEnumW(RESOURCE_GLOBALNET, RESOURCETYPE_ANY, 
                  RESOURCEUSAGE_CONTAINER | RESOURCEUSAGE_CONNECTABLE,
                  nullptr, &hEnum) == NO_ERROR && hEnum) {
    NETRESOURCEW* pBuffer = (NETRESOURCEW*)malloc(16384);
    DWORD count = 0xFFFFFFFF;
    while (WNetEnumResourceW(hEnum, &count, pBuffer, nullptr) == NO_ERROR && count > 0) {
        for (DWORD i = 0; i < count; i++) {
            if (pBuffer[i].lpRemoteName) {
                std::wstring name = pBuffer[i].lpRemoteName;
                if (name.length() >= 2 && name[0] == L'\\' && name[1] == L'\\' &&
                    name.find(L'\\', 2) == std::wstring::npos) {
                    devicesSet.insert(name);
                }
                if (pBuffer[i].dwUsage & RESOURCEUSAGE_CONTAINER) {
                    EnumerateNetworkResourcesRecursive(&pBuffer[i], devicesSet, 0);
                }
            }
        }
        count = 0xFFFFFFFF;
    }
    free(pBuffer);
    WNetCloseEnum(hEnum);
}
\end{lstlisting}

Седьмой метод -- добавление устройств из ARP-таблицы \cite{iphlpapiDOC}, что помогает найти устройства, которые могут отсутствовать в сетевом перечислении.

\begin{lstlisting}[style=CodeListing]
std::vector<std::wstring> arpDevices = GetNetworkIPsFromARP();
for (const auto& dev : arpDevices) {
    if (!dev.empty()) {
        devicesSet.insert(dev);
    }
}
\end{lstlisting}

После применения всех методов множество устройств преобразуется в вектор и возвращается как результат. Если в процессе работы алгоритма произошла ошибка (например, сетевые ресурсы недоступны или доступ запрещён), функция продолжает работу с другими методами, обеспечивая максимально полный список устройств.

\subsubsection{Получение списка сетевых шаров}

Получение списка сетевых шаров на сервере начинается с инициализации вектора для хранения шаров и создания структуры NETRESOURCEW с параметрами для поиска дисковых ресурсов.

Структура NETRESOURCEW настраивается для поиска подключаемых дисковых ресурсов (RESOURCETYPE\_DISK, RESOURCEUSAGE\_CONNECTABLE) на указанном сетевом сервере.

После настройки структуры открывается перечисление сетевых ресурсов через функцию WNetOpenEnumW с указанием пути к серверу. При успешном открытии перечисления выделяется буфер для хранения найденных ресурсов.

\begin{lstlisting}[style=CodeListing]
std::vector<std::wstring> EnumerateNetworkShares(const std::wstring& serverPath) {
    std::vector<std::wstring> shares;
    NETRESOURCEW nr{};
    nr.dwScope = RESOURCE_GLOBALNET;
    nr.dwType = RESOURCETYPE_DISK;
    nr.dwUsage = RESOURCEUSAGE_CONNECTABLE;
    nr.lpRemoteName = const_cast<wchar_t*>(serverPath.c_str());
    
    HANDLE hEnum = nullptr;
    if (WNetOpenEnumW(RESOURCE_GLOBALNET, RESOURCETYPE_DISK, 0, &nr, &hEnum) == NO_ERROR && hEnum) {
        NETRESOURCEW* pBuffer = (NETRESOURCEW*)malloc(16384);
        DWORD count = 0xFFFFFFFF;
        while (WNetEnumResourceW(hEnum, &count, pBuffer, nullptr) == NO_ERROR && count > 0) {
            for (DWORD i = 0; i < count; ++i) {
                if (pBuffer[i].lpRemoteName) {
                    shares.push_back(pBuffer[i].lpRemoteName);
                }
            }
            count = 0xFFFFFFFF;
        }
        free(pBuffer);
        WNetCloseEnum(hEnum);
    }
    return shares;
}
\end{lstlisting}

Выполняется цикл перечисления ресурсов: вызывается функция WNetEnumResourceW для получения списка ресурсов, при успешном выполнении выполняется перебор всех найденных ресурсов, и имена удалённых ресурсов (lpRemoteName) добавляются в вектор шаров.

Цикл продолжается до тех пор, пока не будет получен код ошибки ERROR\_NO\_MORE\_ITEMS, что означает, что все ресурсы перечислены.

После завершения перечисления буфер освобождается, перечисление закрывается через WNetCloseEnum, и вектор шаров возвращается как результат.

Если в процессе работы алгоритма произошла ошибка (например, сервер недоступен, доступ запрещён или перечисление не удалось открыть), функция возвращает пустой вектор шаров.

\subsubsection{Поиск файлов и шаров на сетевых устройствах}

Поиск файлов и шаров на сетевых устройствах начинается с удаления placeholder-элементов из узла сетевого устройства в дереве каталогов.

После удаления placeholder-элементов функция проверяет данные узла на существование и проверяет, был ли узел уже расширен ранее, чтобы избежать повторной загрузки сетевых ресурсов.

Если узел ещё не расширен, проверяется формат сетевого пути. В этом случае вызывается функция EnumerateNetworkShares для получения списка всех доступных шаров на этом сервере.

\begin{lstlisting}[style=CodeListing]
void ExpandNetworkDevice(HWND hTree, HTREEITEM hItem) {
    TVITEMW tvi{}; tvi.mask = TVIF_PARAM; tvi.hItem = hItem; 
    TreeView_GetItem(hTree, &tvi);
    Filesystem::NodeData *data = reinterpret_cast<Filesystem::NodeData*>(tvi.lParam);
    if (!data || data->expanded) return;
    
    if (data->path.length() >= 2 && data->path[0] == L'\\' && 
        data->path[1] == L'\\' && data->path.find(L'\\', 2) == std::wstring::npos) {
        std::vector<std::wstring> shares = EnumerateNetworkShares(data->path);
        for (const auto& share : shares) {
            size_t lastSlash = share.find_last_of(L'\\');
            std::wstring displayName = (lastSlash != std::wstring::npos) ? 
                share.substr(lastSlash + 1) : share;
            bool hasChildren = Filesystem::HasSubdirs(share);
            Filesystem::InsertTreeItem(hTree, hItem, displayName, share, hasChildren);
        }
        data->expanded = true;
    } else {
        Filesystem::ExpandTreeItem(hTree, hItem);
    }
}
\end{lstlisting}

Для каждого найденного шара извлекается отображаемое имя (последний сегмент пути после последней обратной косой черты), проверяется наличие подкаталогов в шаре, и шар добавляется в дерево с соответствующей пометкой о наличии дочерних элементов.

Если путь содержит третью обратную косую черту, это означает, что путь указывает на конкретную шару или подкаталог, и для его обработки вызывается стандартный алгоритм расширения локального узла дерева через функцию ExpandTreeItem.

После завершения обработки узел помечается как расширенный, чтобы предотвратить повторную загрузку при следующем раскрытии.

Если в процессе работы алгоритма произошла ошибка (например, сетевой сервер недоступен, доступ запрещён или шары не найдены), функция завершается без добавления элементов, и узел остаётся в исходном состоянии.